\documentclass[11pt]{article}
\usepackage[a4paper, margin=1in]{geometry}
\usepackage[T1]{fontenc}
\usepackage[utf8]{inputenc}
\usepackage{amsmath, amssymb}
\usepackage{microtype}
\usepackage{parskip}
\usepackage{hyperref}
\hypersetup{colorlinks=true, urlcolor=blue}
\pagestyle{empty}

\begin{document}

\noindent
Claes Johnson, Professor emeritus of Applied Mathematics \\
KTH Royal Institute of Technology, Stockholm \quad\textbar\quad \texttt{claesjohnson@gmail.com}

\vspace{1em}

\noindent
To the Editors, Annales de la Fondation Louis de Broglie \quad\textbar\quad \texttt{inst.louisdebroglie@free.fr} \hfill \today

\vspace{1em}

\noindent Dear Editors,

\medskip

I submit for your consideration the article \emph{Magnetism in RealQM: Real
Charge Currents, Diamagnetism, and the Geometry of Spin}. It is a companion to
my earlier submission on RealQM --- quantum mechanics formulated as $N$ real
charge densities on non-overlapping domains in ordinary three-dimensional
space, in the de Broglie--Schr\"odinger tradition --- and it applies that
real-space ontology to the one arena where the reality of the wave function is
most sharply at stake: magnetism and spin.

The article makes a graded, and deliberately honest, case. \emph{Orbital
magnetism} is shown to be real continuum mechanics: a magnetic moment is a real
circulating current $\mathbf{j}=\rho\mathbf{v}$ carrying real angular momentum,
and the complex wave function $\psi=\sqrt{\rho}\,e^{iS}$ merely bundles the two
real fields $(\rho,\mathbf{v})$ that continuum mechanics has always carried
without any recourse to complex numbers; the sole irreducibly quantum input is
the \emph{integer} winding number, fixed by single-valuedness exactly as
Onsager and Feynman fixed the quantisation of superfluid vortices. \emph{
Diamagnetism} then furnishes a quantitative test: the Langevin susceptibility
$\chi\propto-\langle r^2\rangle$ is read directly from the RealQM density, and a
radial self-consistent field reproduces the measured molar susceptibility of
helium to $0.99\times$, with neon and argon correct in order and trend and the
shortfall diagnosed.

\emph{Spin} is treated as the single residue it is, and I have tried to be
scrupulous about what is and is not established. The article argues that spin
need not be posited as a primitive in RealQM: its exclusion role is already
geometry (non-overlapping domains), its mechanical angular momentum is already
the real circulation of charge (so the Einstein--de Haas effect is native
rather than an obstacle), and the quantisation of its magnitude is the same
topological closure that quantises the orbital moment --- applied to a spinorial
orientation, whose $4\pi$ (SU(2)) periodicity yields $g=2$. By Lévy-Leblond's
theorem this factor is non-relativistic geometry, not a relativistic import, and
it is delivered for one electron by a $\boldsymbol\sigma\cdot\mathbf{p}$
structure with $g=2$ emerging rather than inserted. I state equally plainly what
remains open: that the half-integer closure is exhibited but not yet forced by
the domain geometry, and that closed-shell pairing is \emph{not} supplied by the
spin-blind non-overlap geometry and requires an added exchange term --- a genuine
limitation, reported as such.

I submit to the \emph{Annales} because this is foundational work in the de
Broglie--Schr\"odinger spirit: it asks how much of magnetism follows from taking
the wave function to be a real charge density in three-dimensional space, gives
a quantitative and testable answer, and marks its own boundary honestly rather
than overclaiming. The numerical results are reproducible; the underlying
framework runs in a browser and is documented in a public Gallery.

\medskip

\noindent Thank you for your consideration.

\vspace{1em}

\noindent Sincerely, \\
Claes Johnson

\end{document}
